%*******************************************************************************
%*********************************** First Chapter *****************************
%*******************************************************************************

\chapter{Conclusions and Future Work}  %Title of the First Chapter

\ifpdf
    \graphicspath{{Chapter1/Figs/Raster/}{Chapter1/Figs/PDF/}{Chapter1/Figs/}}
\else
    \graphicspath{{Chapter1/Figs/Vector/}{Chapter1/Figs/}}
\fi


\section{Conclusions}
In this work, NMF based approaches to dereverberation has been studied and tested. An attempt has been made to survey and reproduce results found in literature and tweak the parameters to gain a better understanding of why certain techniques work better. 
First, the NMFD approach, which models reverberation using sub-band filtering works well for dereverberation. The filter, $H$ captures the relation between the source magnitude spectrogram and its reverberated version for each frequency sub-band. 
Also, it is clear that  better model for speech or the room impulse response will definitely help in dereverberation. This is verified by the fact that a better representation for speech, the NMF representation improves the performance of NMFD. 
Convolutive NMF model for speech, although intuitively is a good model for speech, does not improve performance of the dereverberation algorithm. In general, for unsupervised approaches to dereverberation or separation, the convolutive model does not outperform the regular NMF model. It can be concluded that having a more detailed set of basis is only beneficial to capture more information about about the speech and thus is only useful where a supervised approach, using a dictionary of basis is used.
Additionally, different models for the obtained sub-band filter H have been adopted by different authors, but a general consensus on what the room impulse response or H should look like is lacking in literature. Its relation with the magnitude spectrogram of the time domain room impulse response is also unclear and is a possible area to explore. 

\section{Future work}
First, speech recognition rates for the listed approach using a standard recognition system need to be obtained. Since the results obtained do not match quantitative figures quoted in literature and it is difficult to reproduce the experiments exactly, word error rate improvement can provide a benchmark as to where these techniques stand in comparison to state of the art dereverberation techniques. Models for $H$, the set of sub-band envelops, should be explored to understand and improve the assumptions made about it. 
Activation deconvolution technique that was developed does not work as well as techniques in literature, but it does provide an alternative approach to dereverberation. The effects of reverberation on the NMF activation matrix need to be understood and modeled  for this approach to work better.




